\documentclass[10pt,a4paper]{amsart}
\usepackage[margin=19mm]{geometry}
\usepackage{amsmath,amssymb,mathtools}
\usepackage[scr=rsfs,cal=euler]{mathalfa}
\usepackage{xfrac}
\usepackage[unicode,hidelinks,bookmarks=false]{hyperref}
\newcommand{\ie}{i.e.\ }
\newcommand{\cf}{cf.\ }
\newcommand{\resp}{respectively\ }
\newcommand{\Subsection}{\subsection}
\providecommand{\nobreakdash}{\nobreak\hbox{-}}
\setlength{\emergencystretch}{2em}
\multlinegap0pt
\usepackage{bm}
\usepackage{bbm}
\usepackage{dsfont}
\usepackage{xspace}
\usepackage{cancel}
\usepackage{mathtools}
\usepackage{amsthm}
\makeatletter
\@ifundefined{lemma}{\newtheorem{lemma}{Lemma}}{}
\@ifundefined{remark}{\newtheorem{remark}{Remark}}{}
\makeatother
\def\bc{{\mathbb{C}}}
\def\bn{{\mathbb{N}}}
\def\br{{\mathbb{R}}}
\def\bt{{\mathbb{T}}}
\def\bz{{\mathbb{Z}}}
\def\t2deg{\mathbb T^2\text{\rm -deg}}
\def\wt{\widetilde}
\title[Global Bifurcation of Spiral Wave Solutions to the Complex G]{Global Bifurcation of Spiral Wave Solutions to the Complex Ginzburg-Landau Equation}
\author{Carlos Garcia-Azpeitia, Ziad Ghanem, Wieslaw Krawcewicz}
\date{Source: arXiv:2507.15098v2 (CC BY 4.0)}
\begin{document}
\maketitle

\noindent\emph{Source and licence.} Global Bifurcation of Spiral Wave Solutions to the Complex Ginzburg-Landau Equation. Version arXiv:2507.15098v2 (CC BY 4.0), distributed under the Creative Commons Attribution 4.0 International licence, as recorded in the arXiv licence metadata for this exact version and shown on the abstract page https://arxiv.org/abs/2507.15098v2. Full rights notice: licenses/arxiv-2507.15098-CC-BY-4.0.txt.\par\smallskip

\noindent\textit{Editorial scope.} Counted unit: the Lemma whose source label is \texttt{lemm:local\_bif\_inv\_comp\_formula} in the article (source0.tex lines 802--807, source label \texttt{lemm:local\_bif\_inv\_comp\_formula}), delivered together with the article's own complete original proof of it (source0.tex lines 808--829, one \texttt{proof} environment of 22 lines). No mathematics is written, repaired or re-derived: statement and proof are the article's own text line for line. Required context retained uncounted: eq:operator\_equation (lines 554--556); rm:critical\_values (lines 740--765); def:critical\_point\_identification (lines 758--760); eq:local\_bif\_inv\_reformulation (lines 782--787). Every definition, display and label of those lines is retained unchanged. Omitted from this package and not redistributed: 1--553, 557--739, 761--781, 788--801, 830--1050. Nothing inside a retained excerpt is omitted or altered: every retained body is delivered exactly as the article's lines are written, so no omission receipt of any kind accompanies this package. Citations: the retained excerpts carry no citation command at all, so no bibliography entry of the article is transcribed and no reference is redistributed. Reference adaptation: none was needed and none was made; every reference inside the delivered unit resolves inside it, so no unresolved reference or citation is produced. Printed slips kept literal: no printed word, symbol, hypothesis or number of the counted statement or of its proof was altered. Layout and class: the article's own class is \texttt{amsart} with packages including algorithm, algpseudocode, epigraph, mathtools, latexsym, mathrsfs, pstricks, listings, hyperref, pgfplots, appendix, pst-all, pst-plot, mathdots; the wrapper is the lane's pinned \texttt{amsart} editorial wrapper at 10pt on A4, which restates the theorem-like environments the retained text needs and adds the article's own macro definitions that the retained text needs (\texttt{\textbackslash bc}, \texttt{\textbackslash bn}, \texttt{\textbackslash br}, \texttt{\textbackslash bt}, \texttt{\textbackslash bz}, \texttt{\textbackslash t}, \texttt{\textbackslash wt}), with extra wrapper packages (bm, bbm, dsfont, xspace, cancel, mathtools). The delivered numbering is the unit's own. Assets: no retained span contains a figure environment, a graphics-inclusion command, a PDF-page inclusion, a table or a TikZ picture; no image file is copied into this package, the package declares no asset dependency and no image file is redistributed with it. Licence: version arXiv:2507.15098v2 of this article is distributed under the Creative Commons Attribution 4.0 International licence, as recorded in the arXiv licence metadata for this exact version and shown on the abstract page https://arxiv.org/abs/2507.15098v2. A full rights notice is delivered at licenses/arxiv-2507.15098-CC-BY-4.0.txt. Characters: every retained excerpt is pure ASCII, so the delivered engine, XeLaTeX with its default Latin Modern text font, reports no missing character.
\begin{align}\label{eq:operator_equation}
    \mathscr F(\alpha,\beta,u) = 0,
\end{align}

\begin{remark} \label{rm:critical_values} \rm
Recall that a trivial solution $(\alpha_0,\beta_0,0) \in \br \times \br \times \mathscr H$ belongs to the critical set of \eqref{eq:operator_equation} if and only if $0 \in \sigma(\mathscr A(\alpha_0,\beta_0))$. In other words, $(\alpha_0,\beta_0,0) \in \Lambda$ if and only if 
there exist $(m,n) \in \bz \times \bn$
for which 
\begin{align*}
\begin{cases}
\alpha_0 = -s_{|m|,n}; \\
\beta_0 = \omega m  -\eta s_{|m|,n}.
\end{cases}
\end{align*}
Notice that each eigenvalue $\mu_{m,n}: \bc \rightarrow \bc$ admits exactly one root since the zeros of $J'_{|m|}: \br \rightarrow \br$ form a strictly increasing sequence $s_{|m|,0} < s_{|m|,1} < \cdots < s_{|m|,n} < \cdots$. On the other hand, two eigenvalues $\mu_{m,n}$ and $\mu_{m',n'}$ share the same root if and only if
\begin{align*}
 \begin{cases}
s_{|m|,n} = s_{|m'|,n'}; \\
\omega m  - \eta s_{|m|,n} = \omega m'  - \eta s_{|m'|,n'},
\end{cases}   
\end{align*}
which holds if and only if $m = m'$ and $s_{|m|,n} = s_{|m|,n'}$. Again, by strict monotonicity of the sequence  $\{s_{|m|,n}\}_{n \in \mathbb N }$, notice that the latter condition implies $n = n'$. Therefore, each critical point can be uniquely associated with an index pair $(m,n) \in \bz \times \bn$ using the notation 
\begin{align}\label{def:critical_point_identification}
\lambda_{m,n} := \alpha_{m,n} + i \beta_{m,n}, \quad  (\alpha_{m,n}, \beta_{m,n}) :=  \left(-s_{|m|,n}, \; \omega m  - \eta s_{|m|,n} \right),
\end{align}
such that the critical set, in addition to being discrete, admits the following explicit description
\[
\Lambda = \{ (\alpha_{m,n},\beta_{m,n},0) : m \in \bz, \;  n \in  \bn  \}.
\]
\end{remark}

\begin{align}
\label{eq:local_bif_inv_reformulation}
\omega_{\bt^2}(\lambda_0) &= \t2deg\left( 
\bigoplus_{m \in \bz} \bigoplus_{n \in \bn}  \wt{\mathscr A}_{m,n}, \bigoplus_{m \in \bz} \bigoplus_{n \in \bn} \wt{\mathscr O}_{m,n} \right) \\ &= 
\sum_{m \in \bz} \sum_{n \in \bn} \t2deg(\wt{\mathscr A}_{m,n},  \wt{\mathscr O}_{m,n}). \nonumber
\end{align}

\begin{lemma} \label{lemm:local_bif_inv_comp_formula} \rm
Using the notation \eqref{def:critical_point_identification}, the local bifurcation invariant at any critical point $(\lambda_{m,n},0) = (\alpha_{m,n},\beta_{m,n}) \in \Lambda$ is given by the rule
\begin{align}
    \omega_{\bt^2}(\lambda_{m,n}) = (H_m).
\end{align}
\end{lemma}

\begin{proof}
Since each critical point is isotopically simple, i.e. since one has $\mu_{m,n}^{-1}(0) = \{ (\alpha_{m,n},\beta_{m,n})\}$ for all $(m,n) \in \bz \times \bn$ (cf. Remark \ref{rm:critical_values}), the computational formula \eqref{eq:local_bif_inv_reformulation} simplifies to
\begin{align*}
    \omega_{\bt^2}(\lambda_{m,n}) = \t2deg( \wt{\mathscr A}_{m,n}, \wt{\mathscr O}_{m,n}).
\end{align*}
At this point the result follows from the observation that, for all $(\alpha,\beta) \in \br \times \br$, the Jacobian matrix
\begin{align} \label{eq:Jacobian_mu_mnj}
 D \mu_{m,n}(\alpha,\beta) = 
 \frac{1}{(1+i\eta)(s_{|m|,n}+1)}\begin{pmatrix*}[r]
1 & 0 \\  0 & 1
\end{pmatrix*}, \end{align}
is non-singular such that the local Brouwer degrees $\rho_{m,n}(\alpha,\beta)$ are always well-defined and can be computed as follows
\begin{align} \label{eq:rho_mn_formula}
    \rho_{m,n}(\alpha,\beta)
    &= \deg( \det\nolimits_\bc \mathscr A_{m,n},B_{\varepsilon}(\lambda) ) \\ &= \deg(\mu_{m,n}, B_{\varepsilon}(\lambda)) = 
\operatorname{sign}\left[ \det  D \mu_{m,n}(\alpha,\beta)\right]  \nonumber \\ 
 &= \operatorname{sign}\left[  \det
 \begin{pmatrix*}[r]
1 & 0 \\ \nonumber  0 & 1
\end{pmatrix*}\right] = 1. \nonumber
\end{align}
\end{proof}

\end{document}
